\subsection{E4: Robustness to Client Dropout (Simulated)}
\label{sec:e4}

We simulate dropout by withholding a random fraction $p \in \{0, 10, 30, 50\}$\% of clients from each round, under IID and Dirichlet $\alpha{=}0.1$ partitions of the Crop Recommendation dataset; FedAvg aggregates whichever updates arrive. This isolates the learning effect of partial participation; real network disconnections, reconnection and rejoining are measured separately on hardware in E8 and E11. Every configuration is replicated over $n{=}5$ seeds (Table~\ref{tab:e4}).

\begin{table*}[t]
\centering
\caption{E4: Robustness to simulated dropout (Crop Recommendation, AgriMLP, $K{=}10$, $R{=}100$, $n{=}5$ seeds; accuracy as mean $\pm$ half-width of the 95\% CI, in pp). Rounds$\to$40\% is the median across seeds of the first round exceeding 40\% test accuracy. A centralized reference ($K{=}1$, full training set, same loop and seeds) reaches 97.9\% $\pm$ 0.3.}
\label{tab:e4}
\footnotesize
\begin{tabular}{l l r r r r}
\toprule
\textbf{Partition} & \textbf{Dropout} & \textbf{Acc (final)} & \textbf{Acc (peak)} & \textbf{Rounds$\to$40\%} & \textbf{Completed} \\
\midrule
IID & 0\%  & 87.3\% $\pm$ 2.0 & 87.6\% $\pm$ 1.7 & 23 & 100/100 \\
IID & 10\% & 86.5\% $\pm$ 3.4 & 87.4\% $\pm$ 2.8 & 22 & 100/100 \\
IID & 30\% & 86.2\% $\pm$ 3.3 & 88.0\% $\pm$ 2.1 & 22 & 100/100 \\
IID & 50\% & 86.4\% $\pm$ 2.2 & 87.1\% $\pm$ 1.8 & 23 & 100/100 \\
\midrule
Dirichlet $\alpha{=}0.1$ & 0\%  & 71.8\% $\pm$ 5.4 & 72.0\% $\pm$ 5.5 & 33 & 100/100 \\
Dirichlet $\alpha{=}0.1$ & 10\% & 70.7\% $\pm$ 2.5 & 72.4\% $\pm$ 5.0 & 34 & 100/100 \\
Dirichlet $\alpha{=}0.1$ & 30\% & 71.1\% $\pm$ 8.1 & 73.9\% $\pm$ 5.9 & 35 & 100/100 \\
Dirichlet $\alpha{=}0.1$ & 50\% & 67.9\% $\pm$ 6.8 & 69.7\% $\pm$ 5.6 & 37 & 100/100 \\
\bottomrule
\end{tabular}
\end{table*}

Under IID partitioning, accuracy is statistically flat across dropout rates (86.2--87.3\% final, overlapping CIs). Under Dirichlet $\alpha{=}0.1$, label skew costs roughly 15pp, convergence is slower (median 33--37 rounds to 40\%, vs.\ 22--23 under IID), and 50\% dropout shows the lowest mean (67.9\%); with CIs of $\pm$5--8pp no pairwise difference is resolvable at $n{=}5$. Dropout caused neither divergence nor round failures. The federated IID configurations sit $\approx$11pp below the centralized reference (97.9\%), a gap attributable to federation ($K{=}10$, $E{=}1$) rather than to dropout.

\subsection{E5: CoAP Discovery Scaling}
\label{sec:e5}

We compare the two discovery modes of \S\ref{sec:coap} for $N \in \{10, 100, 500, 1{,}000, 2{,}000, 5{,}000\}$ endpoints. Each endpoint carries a domain tag (\texttt{agri}/\texttt{air}), one of ten region tags and two or three sensor tags; queries with one, two and three tags select about 50\%, 5\% and 1--2\% of the nodes. In \emph{probe} mode every endpoint is a separate CoAP server on its own port (hosted by worker processes on one machine, up to $N{=}1{,}000$); in \emph{RD} mode the aggregator runs in a separate process and all endpoints first register concurrently (up to 64 registrations in flight). Wire bytes are the encoded CoAP messages in both directions. Results are medians over 30 lookups (10 for probe) on loopback, so they measure processing and protocol cost; on a WAN, probe cost grows with the RTT of each request, whereas an RD lookup is a single exchange.

\begin{table}[t]
\centering
\caption{E5: discovery latency and wire bytes (median; loopback). Probe = one query per endpoint with client-side filtering; RD = one lookup with server-side filtering.}
\label{tab:e5}
\scriptsize
\setlength{\tabcolsep}{2.5pt}
\begin{tabular}{@{}r l r r r r@{}}
\toprule
\textbf{N} & \textbf{Mode} & \textbf{Tags} & \textbf{Match} & \textbf{ms} & \textbf{Bytes} \\
\midrule
100   & probe & 3 & 1    & 57.0   & 45{,}927 \\
100   & RD    & 3 & 1    & 1.5    & 336 \\
1{,}000 & probe & 1 & 500 & 978.0  & 583{,}649 \\
1{,}000 & probe & 3 & 11  & 795.9  & 459{,}790 \\
1{,}000 & RD    & 1 & 500 & 92.8   & 137{,}519 \\
1{,}000 & RD    & 3 & 11  & 3.1    & 3{,}156 \\
5{,}000 & RD    & 1 & 2{,}500 & 873.7 & 692{,}651 \\
5{,}000 & RD    & 2 & 222 & 66.8   & 61{,}633 \\
5{,}000 & RD    & 3 & 68  & 18.9   & 19{,}626 \\
\midrule
\multicolumn{6}{@{}l}{RD registration, $N{=}5{,}000$: all registered; median 220~ms, p95 329~ms} \\
\bottomrule
\end{tabular}
\end{table}

Probe discovery grows linearly with the number of endpoints regardless of how many match (0.8--1.0~s and 460--580~KB at 1{,}000 nodes), because every node must be queried. With the resource directory, cost scales with the number of \emph{matching} endpoints: a three-tag query over 5{,}000 registered endpoints takes 18.9~ms and 19.6~KB, and paging (\texttt{count}, \texttt{page}) bounds the response size of broad queries. Concurrent registration of all 5{,}000 endpoints completed without failures (median per-registration latency 220~ms with up to 64 in flight). This replaces the original E5, whose loop queried a single server repeatedly and extrapolated byte counts; the corrected probe numbers at 100 nodes (57--97~ms) are of the same order as originally reported.

\subsection{E6: Air Quality Federated Learning}
\label{sec:e6}

\textbf{Dataset and task.} We use the Castilla y Le\'on daily air quality archive~\cite{cyl2023airquality} (10 monitoring zones) merged with AEMET~\cite{aemet2023} meteorological data (wind speed, precipitation), restricted to 2011--2019 to avoid the COVID-19 lockdown artifact~\cite{lopezblanco2024airquality}. The task is 7-day-ahead ICA forecasting: given outdoor readings at day~$t$, predict the ICA class at $t{+}7$. Labels use training-set NO$_2$ terciles ($T_1{=}7.5$, $T_2{=}13.0~\mu$g/m$^3$), balanced in-sample (\textit{Bajo}, \textit{Medio}, \textit{Alto}); because NO$_2$ declined over the decade, the 2019 test year is imbalanced (49.3\% \textit{Bajo}, 27.8\% \textit{Medio}, 22.9\% \textit{Alto}). We therefore report, besides accuracy, macro-F1, balanced accuracy, per-class recall and the confusion matrix (Table~\ref{tab:e6}).

\textbf{Configurations.} AirMLP and AirLSTM (a 7-day input window) with FedAvg under geographic and IID partitions, and ProphetWrapper with FedAvg or FedGAM; $n{=}5$ seeds, $R{=}50$, real Iroh transport. Prophet models are evaluated per zone: each zone's model (federated seasonality plus its local trend under FedGAM; retrained from the full global state under FedAvg) predicts its own test year, and predictions are pooled across zones.

\begin{table*}[t]
\centering
\caption{E6: 7-day-ahead ICA forecasting on the 2019 test year ($K{=}10$, $R{=}50$, $n{=}5$ seeds; mean $\pm$ half-width of the 95\% CI, in \%; real Iroh transport). Class distribution 49.3/27.8/22.9\% (\textit{Bajo/Medio/Alto}).}
\label{tab:e6}
\scriptsize
\setlength{\tabcolsep}{4pt}
\begin{tabular}{l l r r r r r r}
\toprule
\textbf{Model} & \textbf{Partition} & \textbf{Accuracy} & \textbf{Balanced acc.} & \textbf{Macro-F1} & \textbf{Recall \textit{Bajo}} & \textbf{Recall \textit{Medio}} & \textbf{Recall \textit{Alto}} \\
\midrule
Majority class (train)          & ---        & 22.94\% & 33.33\% & 12.44\% & 0.0\% & 0.0\% & 100.0\% \\
Persistence ($t{-}7$)           & ---        & 47.54\% & 43.07\% & 43.13\% & 62.2\% & 30.7\% & 36.3\% \\
\midrule
AirMLP (FedAvg)                 & Geographic & 46.85 $\pm$ 1.02 & 42.22 $\pm$ 0.48 & 40.25 $\pm$ 0.23 & 64.8 $\pm$ 3.7 & 15.3 $\pm$ 2.6 & 46.6 $\pm$ 3.6 \\
AirMLP (FedAvg)                 & IID        & 49.64 $\pm$ 0.61 & 44.93 $\pm$ 0.61 & 43.05 $\pm$ 0.93 & 68.0 $\pm$ 2.0 & 17.2 $\pm$ 2.4 & 49.6 $\pm$ 1.8 \\
AirLSTM (FedAvg)                & Geographic & 53.17 $\pm$ 0.55 & 48.77 $\pm$ 0.42 & 48.11 $\pm$ 0.52 & 69.6 $\pm$ 3.2 & 27.8 $\pm$ 2.9 & 48.9 $\pm$ 2.0 \\
AirLSTM (FedAvg)                & IID        & 55.00 $\pm$ 0.31 & 49.44 $\pm$ 0.34 & 49.07 $\pm$ 0.58 & 74.8 $\pm$ 0.7 & 28.0 $\pm$ 2.5 & 45.6 $\pm$ 1.4 \\
\midrule
ProphetWrapper (FedAvg)         & Geographic & 56.44$^{\ddagger}$ & 48.70 & 49.42 & 79.0 & 36.4 & 30.6 \\
ProphetWrapper (FedAvg)         & IID$^{\S}$ & 57.58 $\pm$ 0.33 & 42.85 $\pm$ 0.35 & 35.63 $\pm$ 0.48 & 63.4 $\pm$ 0.6 & 64.4 $\pm$ 1.2 & 0.8 $\pm$ 0.7 \\
ProphetWrapper (FedGAM)         & Geographic & 56.78$^{\ddagger}$ & 49.18 & 49.94 & 78.9 & 37.5 & 31.1 \\
ProphetWrapper (FedGAM)         & IID$^{\S}$ & 57.93 $\pm$ 0.52 & 42.93 $\pm$ 0.34 & 35.77 $\pm$ 0.47 & 63.9 $\pm$ 0.8 & 64.2 $\pm$ 1.0 & 0.8 $\pm$ 0.7 \\
\bottomrule
\end{tabular}
\vspace{0.3em}

{\scriptsize $^{\ddagger}$Prophet geographic runs are deterministic MAP fits: all five seeds coincide (no CI). $^{\S}$Evaluated on the 2019 test year of a single zone (inherited from the original evaluation code), whose class distribution differs; not comparable with the other rows.}
\end{table*}

\begin{table}[t]
\centering
\caption{E6: confusion matrices summed over five seeds, row-normalized (\%; rows = true class, columns = predicted \textit{Bajo/Medio/Alto})}
\label{tab:e6cm}
\scriptsize
\setlength{\tabcolsep}{3pt}
\begin{tabular}{@{}l l r r r@{}}
\toprule
\textbf{Model (IID)} & \textbf{True} & \textbf{Bajo} & \textbf{Medio} & \textbf{Alto} \\
\midrule
AirMLP  & \textit{Bajo}  & 68.0 & 10.9 & 21.1 \\
        & \textit{Medio} & 47.0 & 17.2 & 35.8 \\
        & \textit{Alto}  & 34.0 & 16.4 & 49.6 \\
\midrule
AirLSTM & \textit{Bajo}  & 74.8 & 14.3 & 10.9 \\
        & \textit{Medio} & 44.5 & 28.0 & 27.6 \\
        & \textit{Alto}  & 28.4 & 26.0 & 45.6 \\
\bottomrule
\end{tabular}
\end{table}

\textbf{Imbalance-aware results.} Accuracy alone overstates the difficulty gap between models and hides how they fail. The 7-day persistence forecaster reaches 47.5\% accuracy but only 43.1\% macro-F1 and balanced accuracy. The memoryless AirMLP does not beat it on the imbalance-aware metrics (macro-F1 40.3\% geographic, 43.1\% IID): it mainly separates the extreme classes and almost never predicts \textit{Medio} correctly (recall 15--17\%), which it confuses with both neighbors (Table~\ref{tab:e6cm}). AirLSTM, which sees a 7-day input window, improves every metric over persistence by about 5--6pp (macro-F1 48.1\% geographic, 49.1\% IID; balanced accuracy 48.8--49.4\%) and nearly doubles the recall of \textit{Medio} (28\%), which remains the hardest class for all models because it is a narrow band between two terciles. The geographic (non-IID) partition costs 1--3pp on every metric for both neural models. The accuracy of AirMLP over real Iroh connections (46.9\% and 49.6\%) matches the values obtained with the in-process transport in the original submission (47.3\% and 49.6\%), a further check that the transport does not alter learning outcomes. The federated Prophet model evaluated over all ten zones (geographic partition) reaches 49.4--49.9\% macro-F1 and 48.7--49.2\% balanced accuracy, on par with AirLSTM and about 6.5pp above persistence; its higher accuracy (56.4--56.8\%) comes largely from the majority class (recall of \textit{Bajo} 79\%, of \textit{Alto} 31\%). The imbalance-aware metrics also exposed a flaw in our original evaluation: in the IID configuration the Prophet models were scored on the test year of a single zone only. On that zone they almost never predict \textit{Alto} (recall below 1\%) and reach only 35.6--35.8\% macro-F1, so the 57.9\% accuracy reported for Prophet-IID in the original submission is not comparable with the other configurations and should not be read as an improvement over persistence. We keep these rows for transparency and base all comparisons on the geographic runs.

\textbf{FedGAM case study.} Switching from FedAvg to FedGAM requires a single server-side change (\texttt{aggregator\_fn}); the same Iroh channels carry AirMLP tensors (1{,}987 parameters) and Prophet Fourier coefficients (27 parameters). Prophet+FedGAM and Prophet+FedAvg are practically indistinguishable (geographic: 49.94\% vs.\ 49.42\% macro-F1, 49.18\% vs.\ 48.70\% balanced accuracy; IID: 35.77\% vs.\ 35.63\% macro-F1). We also tested the regime in which partial aggregation might matter most, with only one year of local history per zone: both rules then fail on this task (geographic macro-F1 18.5\% with FedGAM and 16.7\% with FedAvg, predicting mostly \textit{Alto}), so partial aggregation does not rescue short histories here either. With eight years of local history each zone can estimate its own trend, so restricting aggregation to the seasonal terms brings no measurable benefit here; FedGAM is therefore presented as an example of a model-specific aggregation rule deployed without transport changes, not as a performance contribution.

\subsection{E8: Forced Relay, Outages and Network Switches}
\label{sec:e8}

E8 exercises the relay and the recovery paths on real hardware. A long-lived client on the Raspberry Pi sends a 1~MB update every 5~s to a long-lived server; every transfer re-dials the server by node ID, as the FL data plane does, and the path, active address, retries and local IP addresses are logged per transfer. Perturbations are injected on the Raspberry Pi: blocking all outgoing UDP except DNS with an nftables rule (hole punching becomes impossible and only the relay, which uses HTTPS over TCP, remains), taking the Wi-Fi interface down, and switching between networks with NetworkManager. Throughout E8--E11 a second nftables rule prevents Iroh from using the Tailscale overlay addresses present on the test devices, so that a ``direct'' path is always an Iroh hole-punched path; every transfer is additionally audited for overlay addresses.

\begin{table*}[t]
\centering
\caption{E8: delivery and path per phase (Raspberry Pi on N2 mobile CGNAT $\to$ PC on N1; 1~MB every 5~s; medians)}
\label{tab:e8}
\scriptsize
\setlength{\tabcolsep}{6pt}
\begin{tabular}{@{}l r r r r r@{}}
\toprule
\textbf{Phase} & \textbf{Sent} & \textbf{Delivered} & \textbf{Direct} & \textbf{Relay-carried} & \textbf{Goodput (Mbps)} \\
\midrule
UDP open, start              & 57  & 54 (94.7\%)  & 53  & 1   & \multirow{3}{*}{1.88} \\
UDP open, after unblock 1    & 55  & 40 (72.7\%)  & 39  & 1   & \\
UDP open, after unblock 2    & 57  & 46 (80.7\%)  & 46  & 0   & \\
UDP blocked (2 windows)      & 190 & 190 (100\%)  & 2$^{*}$ & 188 & 1.85 \\
\midrule
\multicolumn{6}{@{}l}{Fallback to relay after blocking: first delivery after 28.1~s and 28.4~s, no transfer lost} \\
\multicolumn{6}{@{}l}{Return to direct after unblocking: 1.9~s and 3.8~s} \\
\bottomrule
\end{tabular}
\vspace{0.3em}

{\scriptsize $^{*}$Classified at the moment of unblocking; one carried a Tailscale address and is excluded from direct counts.}
\end{table*}

\textbf{Forced relay (Table~\ref{tab:e8}).} With UDP blocked, all 190 transfers were delivered over the relay, at the same goodput as the direct path because the 4G uplink, not the relay, was the bottleneck. The switch to the relay took about 28~s, during which the in-flight transfer stalled but was not lost; on unblocking, the path returned to direct within 4~s. On this mobile CGNAT the direct path was less reliable than the relay: 5.3\% of direct transfers were interrupted mid-stream at the start of the run and 19--27\% after the UDP block was lifted, while none of the relay-carried transfers failed. This observation motivated the transfer-level retry of \S\ref{sec:transport}.

\textbf{Outages and network switches.} With transfer retry enabled, a second 30-minute run on N2 delivered 49 of 50 transfers in the first window (two delivered on a retry) and 202 of 213 overall. A 60~s Wi-Fi outage lost no transfer: the in-flight transfer waited and completed, and the first new transfer succeeded 15.6~s after the link came back. During a 120~s outage two transfers timed out, and delivery resumed immediately after the link returned. Switching from 4G to the office Wi-Fi (N2$\to$N1; new local address, new NAT mapping) restored delivery after 4.3~s, now over a direct LAN path to the PC. The opposite switch (N1$\to$N2) did not recover within the remaining five minutes: the Iroh endpoint stopped establishing connections until the process was restarted, even though the 4G link itself was working. We added the endpoint watchdog of \S\ref{sec:transport} as a mitigation (restart with the same key after consecutive failures; covered by the test suite); validating it with repeated network switches in the field remains future work.

\textbf{Long-lived sessions.} Between long-lived nodes the steady-state path is direct: 138 of 140 delivered transfers in the open-UDP phases of the 4G run, and 60 of 60 transfers from the PC on N3$'$ to the Raspberry Pi on N1 (median 266~ms per 64~KB). The exception is a node behind the WSL2 host NAT acting as the \emph{receiver}: in E11 the Raspberry Pi's transfers to and from the aggregator in WSL2 were relay-carried (56 of 58), which we attribute to the additional port-rewriting NAT in front of the server process.

\subsection{E9: Degraded Networks}
\label{sec:e9}

To evaluate behavior on lossy and slow rural links we ran the full FL workload (AirMLP, 10 clients, geographic partition, 20 rounds) and a transfer benchmark (100~KB, 1~MB and 10~MB, 10 transfers each) over real Iroh connections inside an isolated Linux network namespace, applying \texttt{tc-netem} to all traffic: one-way delay of 25--250~ms, jitter, loss of 1--20\%, rate limits of 10~Mbit/s to 512~kbit/s, and two combined profiles. All peers share the impaired link, so a rate limit is shared by concurrent flows. Table~\ref{tab:e9} summarizes the results.

\begin{table*}[t]
\centering
\caption{E9: FL rounds and transfers under emulated network impairments (one-way delay; loss per direction; 20 rounds, 10 clients; medians). Accuracy: final round (best round in parentheses).}
\label{tab:e9}
\footnotesize
\begin{tabular}{l r r r r r r}
\toprule
\textbf{Condition} & \textbf{Rounds OK} & \textbf{Round (s)} & \textbf{Accuracy} & \textbf{Macro-F1} & \textbf{1~MB (s)} & \textbf{10~MB (s)} \\
\midrule
No impairment                          & 20/20 & 1.4  & 46.5\% (47.2) & 38.7\% & 0.13 & 1.2 \\
Delay 25~ms                            & 20/20 & 1.6  & 46.8\% & 39.3\% & 0.60 & 2.2 \\
Delay 100~ms                           & 20/20 & 2.3  & 46.3\% & 38.4\% & 1.86 & 4.6 \\
Delay 250~ms                           & 20/20 & 3.7  & 46.6\% & 38.4\% & 4.37 & 9.5 \\
Delay 100~ms $\pm$ 30~ms               & 20/20 & 2.7  & 46.8\% & 38.1\% & 2.00 & 7.0 \\
Loss 1\%                               & 20/20 & 1.9  & 46.6\% & 40.0\% & 0.15 & 1.5 \\
Loss 5\%                               & 20/20 & 2.5  & 44.3\% & 36.5\% & 0.23 & 1.9 \\
Loss 10\%                              & 20/20 & 3.3  & 48.4\% & 39.1\% & 0.27 & 2.6 \\
Loss 20\%                              & 20/20 & 8.7  & 26.8\% (46.9) & 21.9\% & 1.70 & 10.6 \\
Rate 2~Mbit/s                          & 20/20 & 2.7  & 46.5\% & 38.9\% & 4.49 & 43.8 \\
Rate 512~kbit/s                        & 20/20 & 10.7 & 46.7\% & 38.9\% & 17.0 & 167 \\
Rural LTE (30$\pm$10~ms, 2\%, 5~Mbit/s) & 20/20 & 4.1 & 46.8\% & 38.3\% & 2.14 & 27.2 \\
Harsh rural (75$\pm$25~ms, 10\%, 1~Mbit/s) & 20/20 & 6.5 & 46.6\% & 38.5\% & 26.2 & failed$^{\dagger}$ \\
\bottomrule
\end{tabular}
\vspace{0.3em}

{\scriptsize $^{\dagger}$0/10 transfers of 10~MB completed within the 180~s timeout. A jitter profile that reorders packets (normal-distributed delay without rate shaping) stalled the run and is excluded; reordering of that magnitude is atypical of real links.}
\end{table*}

All 260 rounds completed in every condition. For the 14~KB updates of AirMLP, impairments lengthen rounds (from 1.4~s to 3.7~s at 250~ms delay, 8.7~s at 20\% loss and 10.7~s at 512~kbit/s) but do not change final accuracy or macro-F1, which stay within the seed-to-seed range of the unimpaired run, with one exception: at 20\% loss two rounds aggregated only 8 and 5 of the 10 clients because updates arrived after the round timeout, and the last round was the 5-client one, so its non-IID aggregate drops to 26.8\% (the best round of that run reached 46.9\%). Payload size is what makes QUIC's behavior visible: 10~MB transfers grow from 1.2~s to 9.5~s at 250~ms delay and 167~s at 512~kbit/s, and fail entirely under the harsh rural profile. For edge-sized models the transport is not the limiting factor; for multi-megabyte models on such links, update compression or longer round deadlines would be required.

\subsection{E10: Public-Key Admission Control}
\label{sec:e10}

We measure admission control (\S\ref{sec:admission}) on loopback and over the WAN (aggregator on N3$'$, Raspberry Pi on N1). An authorized client holding an allow-listed key and an unauthorized client with a fresh key each attempt 30 transfers; we also test a peer that registers another node's endpoint over \texttt{fl-ctrl/1} (Table~\ref{tab:e10}).

\begin{table}[t]
\centering
\caption{E10: admission control ($N{=}30$ attempts each; medians)}
\label{tab:e10}
\scriptsize
\setlength{\tabcolsep}{3pt}
\begin{tabular}{@{}l r r@{}}
\toprule
\textbf{Metric} & \textbf{Loopback} & \textbf{WAN (N1$\to$N3$'$)} \\
\midrule
Unauthorized payloads accepted        & 0/30     & 0/30 \\
Unauthorized attempts rejected/logged & 30/30    & 30/30 \\
Time to rejection (client view)       & 10.1~ms  & 329~ms \\
Authorized transfers delivered        & 30/30    & 30/30 \\
Authorized latency, allow-list off/on & 11.0 / 11.1~ms & --- / 3.97~s$^{\ddagger}$ \\
Spoofed registration rejected         & yes      & --- \\
\bottomrule
\end{tabular}
\vspace{0.3em}

{\scriptsize $^{\ddagger}$64~KB, relay-carried (receiver behind the WSL2 host NAT).}
\end{table}

No unauthorized payload was ever accepted, rejection costs one handshake plus one round trip (10~ms on loopback, 0.33~s over the WAN), and the allow-list check adds no measurable latency to authorized transfers. Because the check runs on the authenticated key before any stream is read, an unauthorized peer cannot inject an update; together with the per-round binding of updates to selected node IDs it also prevents a member from submitting several updates per round.

\subsection{E11: Federation over the WAN on a Raspberry Pi~3B}
\label{sec:e11}

Finally, we ran a complete federation across two sites: the aggregator on the PC at home (N3$'$) and three clients registered over \texttt{fl-ctrl/1} by public key---the Raspberry Pi~3B in the office (N1, zone 0) and two clients on the PC (zones 1 and 2)---with the air-quality task, AirMLP, $R{=}30$ and a minimum of two clients per round. No port forwarding, VPN or public IP was configured. The Raspberry Pi recorded CPU time, resident memory and system load per round phase (Table~\ref{tab:e11}). The Raspberry Pi~3B has no power sensor, so energy is estimated with a linear model $P = P_\text{idle} + (P_\text{max}-P_\text{idle})\,u$ on the system CPU utilization $u$, with board-level $P_\text{idle}{=}1.4$~W and $P_\text{max}{=}3.7$~W~\cite{geerlingpower}; these values are estimates, not measurements.

\begin{table}[t]
\centering
\caption{E11: per-round resource use on the Raspberry Pi~3B client (AirMLP, 14~KB updates; medians over 29 rounds; energy estimated)}
\label{tab:e11}
\scriptsize
\setlength{\tabcolsep}{3pt}
\begin{tabular}{@{}l r r r r@{}}
\toprule
\textbf{Phase} & \textbf{Wall (s)} & \textbf{CPU (s)} & \textbf{Power (W)} & \textbf{Energy (J)} \\
\midrule
Receive model (incl.\ waiting) & 3.10 & 0.11 & 1.46 & 4.5 \\
Local training (1 epoch)       & 1.02 & 2.79 & 2.86 & 2.9 \\
Send update                    & 3.45 & 0.20 & 1.51 & 5.2 \\
\midrule
\multicolumn{5}{@{}l}{Resident memory: 255~MB after start-up, 419~MB peak (of 1~GB)} \\
\bottomrule
\end{tabular}
\end{table}

The Raspberry Pi joined the federation shortly after the first round and took part in every subsequent round: 29 of 30 rounds completed, 27 with all three clients, at a median of 7.9~s per round (p95 21.3~s). The last round failed because of a race we found in this experiment: when the receiver closed a connection right after reading a model, the sender sometimes saw the close before the acknowledgement and resent the model, so clients occasionally trained twice per server round and finished early (four stale updates were discarded by the per-round binding of \S\ref{sec:admission}). Receivers now close with an explicit receipt code that senders treat as delivery; with this fix the live comparison of \S\ref{sec:flower}, on the same hardware, completed 30 of 30 rounds with all clients and no discarded update. Local training is cheap on this device---1.0~s of wall time and 2.8~s of CPU time (multi-core) per epoch---while the communication phases are dominated by waiting for other clients and by the relay-carried path to the aggregator in WSL2 (56 of 58 Raspberry Pi transfers; \S\ref{sec:e8}). Memory, not computation, is the tighter constraint: the Python/PyTorch runtime occupies 255--419~MB, about 40\% of the device's RAM, which supports our recommendation of a native client runtime for smaller devices (\S\ref{sec:agnostic}). We also note a deployment pitfall: the default PyPI PyTorch wheel for 64-bit ARM is built for newer ARMv8.2 server cores and terminates with an illegal-instruction error during back-propagation on Cortex-A53/A72 boards; the CPU-only wheel works.

\subsection{Operational Comparison with Flower over Tailscale}
\label{sec:flower}
